%\documentclass[11pt,a4paper,uplatex,dvipdfmx]{ujarticle} 		% for uplatex
\documentclass[11pt,a4paper]{ltjsarticle}
\input{pieces/form00_header} % pieces
\input{pieces/kakenhi7} % pieces
\input{pieces/form01_dcpd_header} % pieces
\input{pieces/hook3} % pieces
%#Name: dc
\input{pieces/form03_dcpd_headers} % pieces
\input{pieces/form04_dc_header} % pieces
% ===== Global definitions for the Kakenhi form ======================
%　基本情報
%
%------ 研究課題名  -------------------------------------------
\newcommand{\研究課題名}{軌道体射影直線の導来圏の半直交分解とBridgeland安定性条件に関する研究}

%----- 研究機関名と研究代表者の氏名-----------------------
\newcommand{\研究機関名}{大阪大学}
\newcommand{\研究代表者氏名}{山本 雄大}
\newcommand{\me}{\underline{\underline{}}}
\input{pieces/inst_general_images} % pieces
% user07_header
% ===== my favorite packages ====================================
% ここに、自分の使いたいパッケージを宣言して下さい。
\usepackage{wrapfig}
\usepackage{multicol,amsmath,amssymb,mathtools,ascmac,amsthm,amscd,physics,comment,dcolumn,mathrsfs,test,tikz-cd,enumitem}
%\usepackage{amssymb}
%\usepackage{mb}
%\DeclareGraphicsRule{.tif}{png}{.png}{`convert #1 `dirname #1`/`basename #1 .tif`.png}
\usepackage{lineno}

% ===== my personal definitions ==================================
% ここに、自分のよく使う記号などを定義して下さい。
\newcommand{\klpionn}{K_L \to \pi^0 \nu \overline{\nu}}
\newcommand{\kppipnn}{K^+ \to \pi^+ \nu \overline{\nu}}

% ----- 業績リスト用 -------------
\newcommand{\paper}[6]{%
	% paper{title}{authors}{journal}{vol}{pages}{year}
	\item ``#1'', #2, #3 {\bf #4}, #5 (#6).			% お好みに合わせて変えてください。
}

\newcommand{\etal}{\textit{et al.\ }}
\newcommand{\ca}[1]{*#1}	% corresponding author;   \ca{\yukawa}  みたいにして使う
\newcommand{\invitedtalk}{招待講演}

\newcommand{\yukawa}{H.~Yukawa}					% no underline
%\newcommand{\yukawa}{\underline{\underline{H.~Yukawa}}}	% with 2 underlines
\newcommand{\tomonaga}{S.~Tomonaga}

\newcommand{\prl}{Phys.\ Rev.\ Lett.\ }		% よく使う雑誌も定義すると楽

% ===== 欄外メモ ==================
\newcommand{\memo}[1]{\marginpar{#1}}
%\renewcommand{\memo}[1]{}	% 全てのメモを表示させないようにするには、行頭の"%"を消す

%\input{../../sample/simple/contents}	% skip
\input{pieces/hook5} % pieces

\begin{document}
\input{pieces/hook7} % pieces
%#Split: 01_background  
%#PieceName: p01_background
\input{pieces/p01_background_00}
\section{研究の概要と位置づけ}
%　　　　＜＜最大　1ページ＞＞

%s03_abst_background
\noindent
\underline{\textbf{研究課題名：\研究課題名}}

%begin 研究目的や背景などの概要 ＝＝＝＝＝＝＝＝＝＝＝＝＝＝＝＝＝＝＝＝
\bullet\ 研究の概要\\ 本研究では，軌道体射影直線の導来圏に対する Bridgeland 安定性条件の空間を具体的に記述することを目的とする．軌道体射影直線は，表現論と代数幾何を結びつける基本的な Deligne--Mumford スタックであり，ADE 型箙の導来圏が半直交分解の成分として現れることから，既知の安定性条件を貼り合わせることで，全体の安定性条件空間を構成するための自然な研究対象となる．Bridgeland 安定性条件は，導来圏の対象を「安定な構成要素」に分解する理論であり，その全体は複素多様体を成す．この空間は，自己同値群の作用，モジュライ空間の壁越え現象，ホモロジー的ミラー対称性と深く関わるため，導来圏に付随する最も重要な幾何学的対象の一つと考えられている．

本研究で扱う軌道体射影直線は，その導来圏には半直交分解や例外対象列が存在する．特に，その部分圏として ADE 型箙の導来圏が現れることが知られている．ADE 型箙の導来圏については，安定性条件空間と配置空間，組紐群作用との関係が既に詳しく研究されているが，軌道体射影直線の導来圏については，部分圏上の安定性条件から全体の安定性条件をどのように構成するかは未解決である．

そこで本研究では，半直交分解の各成分上で既に理解されている安定性条件を比較し，それらを貼り合わせることで，全体の導来圏上の安定性条件を構成する条件を定式化する．さらに，この方法を用いて，安定性条件空間の連結性や自己同値群の作用を具体的に記述し，ADE 型の場合に知られている配置空間やその被覆空間との対応が，軌道体射影直線の場合にどのように現れるかを調べる．これにより，半直交分解という導来圏の離散的構造から，安定性条件空間という連続的な幾何学的構造を復元する方法を確立することを目指す．また，将来的には，半直交分解と各成分上の安定性条件から，全体の導来圏上の Bridgeland 安定性条件を体系的に構成する一般理論への発展を見据えている．\par
\bullet\ 研究の位置づけ\\
Bridgeland 安定性条件は，導来圏の対象を「安定な構成要素」に分解する理論であり，
その全体は複素多様体を成す．
この空間の構造を理解することは，
自己同値群の作用，モジュライ空間の壁越え，
ホモロジー的ミラー対称性を理解する上で重要である．
ADE 型箙の導来圏については，
Bessis による組紐群と配置空間の研究，
Hertling による Frobenius 多様体の研究，
および Haiden--Katzarkov--Kontsevich による有限型・アフィン A 型の場合の研究により，
安定性条件空間の位相的・幾何学的構造が詳しく調べられている．
これらの研究では，
安定性条件空間が配置空間の被覆として記述され，
組紐群や Coxeter 変換との関係が明らかにされている．
一方，軌道体射影直線の導来圏については，
その内部に ADE 型箙の導来圏が半直交分解の成分として現れるにもかかわらず，
部分圏上の安定性条件から全体の安定性条件をどのように構成するかは
未解決である．
特に，半直交分解という導来圏の離散的構造から，
安定性条件空間という連続的な幾何学的構造をどこまで復元できるかは，
導来圏論における基本的問題である．
申請者はこれまでに，
Bridgeland 安定性条件の定義と基本性質，
および ADE 型箙の導来圏における Gepner 型安定性条件について学び，
軌道体射影直線の導来圏との関係を整理してきた．
その過程で，
半直交分解の各成分で理解されている安定性条件を
貼り合わせることで，
全体の導来圏上の安定性条件を構成できるのではないかという着想に至った．
本研究は，
ADE 型で確立された理論を基盤として，
より一般の Deligne--Mumford スタックへ拡張し，
表現論，代数幾何，ミラー対称性を結びつける
新たな枠組みを与える研究として位置づけられる．
%end 研究目的や背景などの概要 ＝＝＝＝＝＝＝＝＝＝＝＝＝＝＝＝＝＝＝＝


%begin 本研究の着想に至った経緯など ＝＝＝＝＝＝＝＝＝＝＝＝＝＝＝＝＝＝＝＝
%end 本研究の着想に至った経緯など ＝＝＝＝＝＝＝＝＝＝＝＝＝＝＝＝＝＝＝＝

\input{pieces/p01_background_01}

%#Split: 02_purpose_plan  
%#PieceName: p02_purpose_plan
\input{pieces/p02_purpose_plan_00}
\section{【２】研究計画（２）研究目的・内容等}\\
\textcircled{1}\ 本研究の目的は，軌道体射影直線の導来圏における Bridgeland 安定性条件の空間を，半直交分解および ADE
型箙の導来圏との関係を通して明らかにすることである．三角圏は代数幾何や表現論に現れる対象を統一的に扱う
基本的な枠組みであるが，一般には対象の分解構造を直接記述することが難しい．Bridgeland 安定性条件は，三
角圏の対象に Harder–Narasimhan 分解を与え，圏の対象を安定な構成要素に分解して理解するための有効な道具
である．また，安定性条件全体の空間は複素多様体となることが Bridgeland によって示され，圏に付随する幾何的パラメータ空間として振る舞う．
そのため，この空間の構造を調べることは，圏の内部構造や自己同値，不変量を理解する上で重要である．

 さらに，ホモロジー的ミラー対称性の観点からは，この安定性条件の空間がミラー多様体（あるいは対応する Landau--Ginzburg 模型）の複素構造変形空間や Frobenius 多様体構造を記述するモジュライ空間と一致することが予想されており，この幾何学的対応を検証する上でも極めて重要である．特に ADE 型特異点やアフィン A 型箙などの先行研究では，安定性条件の空間が特異点の臨界値配置から得られる配置空間の被覆空間として実現されることが知られている．本研究の対象である軌道体射影直線は，その内部に ADE 型箙の導来圏が半直交分解の成分としてとして現れるため，導来圏の半直交分解という代数的構造から，ミラー側で現れる連続的な幾何学的・位相的情報をどこまで正確に読み取れるかという，ミラー対称性理論における問いに対して具体的な解答を与えるための重要な鍵を握っている．\vspace{0.5mm}\\

② そこで本研究では，軌道体射影直線の導来圏が有する半直交分解のデータと，その分解後の各部分圏として
現れる ADE 型箙の導来圏の構造に着目し，構成的に全体を統括するアプローチをとる．これら個々の ADE
型箙の導来圏については，すでに Gepner 型安定性条件の存在や，組紐群作用と配置空間の被覆構造との関係
など，安定性条件空間の性質が先行研究によって深く理解されている．したがって，本研究の第 1 段階として，
これら既知である「半直交分解のデータ」と「各成分上の安定性条件の情報」を網羅的に統合・比較し，
貼り合わせることで，全体の導来圏上の Bridgeland 安定性条件を具体的に構成することから研究を開始する．
特に，軌道体射影直線の導来圏には，
\[
D^b(\mathrm{Coh}\,\mathbb{P}_{2,3,5})
=
\langle \mathcal{E}, D^b(kE_8)\rangle .
\]
のように，ADE 型箙の導来圏が部分圏として現れる半直交分解が存在する．ADE 型の場合には，特異点の臨
界値配置から得られる配置空間の被覆空間と，Bridgeland 安定性条件の空間が対応することが，ミラー対称性の立
場から予想されている．本研究ではまず，軌道体射影直線の導来圏に現れる半直交分解と，その部分圏として現れ
る ADE 型箙の導来圏を用いて，この予想に従い安定性条件空間を具体的に構成することを目指す．特に，各成分
上で理解されている安定性条件を比較し，それらを貼り合わせることで，全体の導来圏上の安定性条件を構成する
条件を定式化する．さらに，将来的には，この構成法をより一般の半直交分解や Deligne–Mumford スタックへ拡
張し，半直交分解と各成分上の安定性条件から，全体の導来圏上の Bridgeland 安定性条件を体系的に構成する一
般理論への発展を目指す．

さらに，この軌道体射影直線における具体的な構成プロセスを足がかりとすることで，研究を第 2 段階へと進展
させる．すなわち，この特定の幾何学的スタックにおける成功例を抽象化し，一般的な三角圏において「半直交
分解」と「分解された各成分上の安定性条件」があらかじめ与えられた際に，それらの代数的・離散的な構造
データのみから全体の導来圏上の安定性条件を体系的に復元できるような，一般理論へと拡張・発展
させることを目指す．これにより，これまで個別の幾何学的対象ごとに依存していた安定性条件の構成法を，
圏の分解という代数的な枠組みから記述する新たな枠組みを確立する．

\begin{center}
\fbox{
\parbox{0.9\textwidth}{
\centering
\textbf{軌道体射影直線の導来圏}

\vspace{1mm}
{\small 半直交分解・例外対象列を通して解析}

\vspace{2mm}
$\Downarrow$

\vspace{2mm}
\textbf{ADE 型箙の導来圏との関係}

\vspace{1mm}
{\small 既知の安定性条件空間・組紐群作用を利用}

\vspace{2mm}
$\Downarrow$

\vspace{2mm}
\textbf{Bridgeland 安定性条件空間の構成}

\vspace{1mm}
{\small 配置空間・被覆空間との対応を検証}
}
}
\end{center}\vspace{3mm} 
\textcircled{3}\ 本研究の独創的な点は，
Bridgeland 安定性条件空間を個別に直接計算・記述するのではなく，
半直交分解という導来圏の離散的・代数的データから，
安定性条件空間という連続的・幾何学的対象を構成的に復元しようとする点にある．\par
本研究が完成すれば，軌道体射影直線の導来圏における Bridgeland 安定性条件空間の構造を，
半直交分解および ADE 型箙の導来圏との関係を通じて，
具体的かつ体系的に記述できると期待される．
Bridgeland 安定性条件空間は，
導来圏に付随する基本的な幾何学的対象であり，
その連結性，基本群，自己同値群の作用などを理解することは，
導来圏の対称性や変形理論を調べる上で重要な課題である．
特に，
スタック的特異点をもつ代数多様体に対して，
導来圏の内部構造からどこまで大域的な幾何学的・位相的情報を復元できるかは，
現在も十分には理解されていない．

本研究では，
半直交分解という導来圏の離散的・代数的な分解データと，
各成分上で既に理解されている安定性条件の情報を貼り合わせることで，
全体の安定性条件空間を構成することを目指す．

さらに，本研究が一般化されれば，
「分解された各成分上で理解されている構造や性質を，この分野に限らず
どのように全体へ拡張・復元できるか」という問題に対する新たな方法論につながる可能性がある．

本研究の中心的な問題意識の一つは，半直交分解という導来圏の「離散的」な構造データから，安定性条件空間の
ような「連続的」な幾何学的対象の情報をどこまで復元できるかという点にある．
本研究によって安定性条件空間の構造が明示的に記述されれば，スタックの導来圏においても，箙の導来圏で知ら
れている自己同値群の作用や組紐群的構造を統一的に理解できるようになる．この成果は，表現論においても重要
な意義をもつ．ADE 型箙の導来圏では，例外対象列の変異，組紐群作用，Coxeter 変換などが豊かな対称性を形成
することが知られているが，これらが安定性条件空間の幾何とどのように対応するかは中心的な課題である．本研
究により，例外対象列の組合せ論，半直交分解の離散的データ，自己同値群の作用，安定性条件空間の位相的・複
素幾何学的構造との関係が具体的に記述され，導来圏に現れる対称性を統一的に理解する枠組みが得られることが
期待される．

%end 研究目的と研究計画shorter ＝＝＝＝＝＝＝＝＝＝＝＝＝＝＝＝＝＝＝＝
\input{pieces/p02_purpose_plan_01}

%#Split: 03_rights  
%#PieceName: p03_rights
\input{pieces/p03_rights_00}
\section{人権の保護及び法令等の遵守への対応}
%　　　　＜＜最大　1ページ＞＞

% s09_rights
%begin 人権の保護及び法令等の遵守への対応 ＝＝＝＝＝＝＝＝＝＝＝＝＝＝＝＝＝＝＝＝
特に該当なし．

%end 人権の保護及び法令等の遵守への対応 ＝＝＝＝＝＝＝＝＝＝＝＝＝＝＝＝＝＝＝＝

\input{pieces/p03_rights_01}

%#Split: 04_abilities  
%#PieceName: p04_abilities
\input{pieces/p04_abilities_00}
\section{【４】研究遂行力の自己分析}
%　　　　＜＜最大　2ページ＞＞

% s14_abilities
%\SelfReviewInstructions\\% <-- 留意事項：これは消すか、コメントアウトしてください。
\noindent
\bullet\ (1) 研究に関する自身の強み\\
私は，個別の対象に対する具体的計算を通して，その背後にある一般構造を捉えることを重視して研究を進め
ている．特に，導来圏に現れる半直交分解や例外対象列のような離散的構造と，Bridgeland 安定性条件空間
のような連続的幾何との関係について，一貫して比較・整理を行っている．現在の研究計画も，「導来圏の内
部構造から，そこに付随する幾何学的・位相的対象をどこまで復元できるのか」という問題意識に基づいてい
る．

また私は，大局的な構造や理論的枠組みを意識しながら研究を進める一方で，具体例に対する実際の計算を
通して理論を検証することを重視している．抽象理論のみを形式的に追うのではなく，具体的対象において何
が起きているのかを丁寧に調べることで，理論の本質や限界を把握しようと努めている．特に，軌道体射影直
線や ADE 型箙の導来圏のように，具体計算と一般理論の双方が現れる対象に着目することで，抽象理論の背
後にある幾何学的構造や対称性を捉えるよう努めている．その中で，個別の計算から共通する構造を見出し，
それを一般化する方向へ研究を発展させる姿勢を意識している．

研究を進める際には，まず既知の具体例に対して理論がどのように振る舞うかを丁寧に整理し，そこから共
通する構造や一般化可能な条件を抽出することを意識している．特に，導来圏論では同じ対象が複数の異なる
方法で記述されることが多いため，それぞれの構成や対応関係を比較しながら理解することを重視している．
現在の研究でも，ADE 型箙の導来圏で知られている安定性条件空間の構造を出発点として，軌道体射影直線
の導来圏に現れる半直交分解との対応を整理し，どの情報が全体の安定性条件空間の構成に本質的であるかを
段階的に検証している．このように，具体例を通して問題点や構造を把握し，そこから一般理論へ接続してい
く研究姿勢を重視している．

また，異なる分野に現れる構造の共通性に強い関心を持っていることも，自身の特徴であると考えている．
現在関心を持っている Bridgeland 安定性条件は，表現論，代数幾何，特異点理論，ミラー対称性など，多く
の分野と深く関係している．特に，組紐群作用，配置空間，被覆空間といった位相的・幾何学的構造が，導来
圏や安定性条件とどのように結びつくかに関心を持っている．私は，個別分野の技術を個別に扱うのではなく，
異なる分野で現れる構造を比較・整理しながら，その背後にある共通原理を捉えることを重視して研究を進め
ている．

さらに，研究を進める際には，長期的な問題意識を持ちながら，段階的に具体例を積み重ねていくことを意
識している．現在の研究計画でも，まず軌道体射影直線と ADE 型の場合に対して安定性条件空間を具体的に
構成し，その後より一般の半直交分解へ拡張するという段階的な方針を立てている．特に，既知の理論が成立
している具体例において，どの条件が本質的であり，どの部分が一般化可能であるかを整理しながら研究を進
めている．このように，大局的な構想を持ちながらも，具体例による検証を通して一歩ずつ理論を構築してい
く姿勢が，自身の研究上の強みであると考えている．

また，自身の研究内容を専門外の研究者にも伝えられるよう，問題意識や背景を意識して説明することを重
視している．導来圏論や安定性条件は抽象度の高い分野であるが，本研究の本質には，「離散的構造から連続
的幾何を復元する」という数学全体に共通する問題意識が存在している．そのため，個別の技術的議論だけで
なく，研究の背後にある数学的意義や，他分野との接点を明確に意識しながら研究を進めている．

\vspace{2mm}

\bullet\ (2) 今後研究者として更なる発展のため必要と考えている要素\\
今後さらに研究を発展させるためには，個別の理論を深く理解するだけではなく，異なる分野を横断して構
造を統一的に捉える視点をさらに養う必要があると考えている．現在の研究テーマである Bridgeland 安定性
条件は，導来圏論に留まらず，特異点理論，位相幾何，数理物理，ミラー対称性など多くの分野と関係してい
る．そのため，将来的には，個々の技術を習得するだけではなく，異なる分野に現れる構造の対応関係をより
深く理解し，新しい問題設定や理論形成につなげられる研究者になりたいと考えている．
%\PapersInstructions\\ %<-- 留意事項：これは消すか、コメントアウトしてください。
%begin 自己分析 ＝＝＝＝＝＝＝＝＝＝＝＝＝＝＝＝＝＝＝＝
\vspace{5mm}
\noindent
%begin 今後必要な要素 ＝＝＝＝＝＝＝＝＝＝＝＝＝＝＝＝＝＝＝＝
%end 今後必要な要素 ＝＝＝＝＝＝＝＝＝＝＝＝＝＝＝＝＝＝＝＝

\input{pieces/p04_abilities_01}

%#Split: 99_tail
\input{pieces/hook9} % pieces
\end{document}

